% This document was generated by an LLM.

\documentclass{essay}

\title{Self-Check Questions:\\Dedekind Cuts}
\author{}
\date{}

\begin{document}

\maketitle
\tableofcontents
\newpage

\section{Definitions of a Dedekind Cut}

\begin{enumerate}[label=\arabic*., itemsep=10pt]

\item State the two-set definition of a Dedekind cut $(A,B)$ of $\mathbb{Q}$.

\begin{answer}
A Dedekind cut is a pair $(A,B)$ of nonempty subsets of $\mathbb{Q}$ such that
\[
A\cup B=\mathbb{Q},\qquad A\cap B=\varnothing,
\]
every $a\in A$ is less than every $b\in B$, and $A$ has no greatest element.
\end{answer}

\item State the one-set definition using only the lower set $A\subseteq\mathbb{Q}$.

\begin{answer}
A Dedekind cut may be defined as a nonempty proper subset $A\subsetneq\mathbb{Q}$ that is downward closed and has no greatest element.
\end{answer}

\item Why is it enough to specify only $A$? How is $B$ recovered?

\begin{answer}
Because $A$ and $B$ partition $\mathbb{Q}$,
\[
B=\mathbb{Q}\setminus A.
\]
Thus either half uniquely determines the other.
\end{answer}

\item What does it mean for $A$ to be downward closed?

\begin{answer}
Downward closed means
\[
x\in A,\ y<x \Longrightarrow y\in A.
\]
\end{answer}

\item Write the statement ``$A$ has no greatest element'' entirely with quantifiers.

\begin{answer}
\[
\forall a\in A\;\exists a'\in A\text{ such that }a<a'.
\]
\end{answer}

\item Are the following two statements equivalent?
\[
\text{$A$ has no greatest element}
\]
and
\[
\forall a\in A\;\exists a'\in A\text{ such that }a<a'.
\]
Explain.

\begin{answer}
Yes. If every $a\in A$ has a strictly larger element of $A$, then no element can be greatest, and conversely.
\end{answer}

\item Why must a Dedekind lower set satisfy both
\[
A\neq\varnothing
\qquad\text{and}\qquad
A\neq\mathbb{Q}?
\]

\begin{answer}
$A\neq\varnothing$ prevents the trivial empty lower half, while $A\neq\mathbb{Q}$ ensures that there is a nonempty upper half.
\end{answer}

\item Suppose $A$ is nonempty, proper, and downward closed, but is allowed to have a greatest element.
What extra objects would then be admitted?

\begin{answer}
Such sets would include versions of rational cuts whose lower half contains its endpoint, for example
\[
\{q\in\mathbb{Q}:q\le 3\}.
\]
This creates two possible cuts for the same rational boundary unless a convention is imposed.
\end{answer}

\item Why is the ``no greatest element'' condition partly a convention about where a rational boundary is placed?

\begin{answer}
For a rational boundary $r$, one must decide whether $r$ belongs to the lower or upper half. Requiring $A$ to have no greatest element forces $r$ into $B$.
\end{answer}

\item Give the cut corresponding to the rational number $3$ using the usual lower-set convention.

\begin{answer}
\[
A=\{q\in\mathbb{Q}:q<3\},\qquad
B=\{q\in\mathbb{Q}:q\ge 3\}.
\]
\end{answer}

\item Could one instead define rational cuts using
\[
A=\{q\in\mathbb{Q}:q\le r\}?
\]
What would have to change in the definition?

\begin{answer}
Yes, but then the axioms would need to be changed, for example by requiring $B$ to have no least element instead. It is a different endpoint convention.
\end{answer}

\item State the order-theoretic definition of a Dedekind cut using the phrase ``initial segment.''

\begin{answer}
A Dedekind cut is a nonempty proper initial segment of $(\mathbb{Q},<)$ with no greatest element.
\end{answer}

\end{enumerate}

\section{Maximum, Upper Bound, and Supremum}

\begin{enumerate}[label=\arabic*., itemsep=10pt]

\item Define what it means for $m$ to be the maximum of a set $S$.

\begin{answer}
$m$ is the maximum of $S$ if $m\in S$ and
\[
x\le m
\]
for every $x\in S$.
\end{answer}

\item Define what it means for $u$ to be an upper bound of $S$.

\begin{answer}
$u$ is an upper bound of $S$ if
\[
x\le u
\]
for every $x\in S$. It need not belong to $S$.
\end{answer}

\item Define what it means for $s$ to be the supremum of $S$.

\begin{answer}
$s$ is the supremum of $S$ if $s$ is an upper bound and every upper bound $u$ satisfies
\[
s\le u.
\]
\end{answer}

\item What is the logical difference between
\[
\max S
\qquad\text{and}\qquad
\sup S?
\]

\begin{answer}
A maximum must belong to $S$. A supremum need not belong to $S$.
\end{answer}

\item Can a set have a supremum without having a maximum? Give an example in $\mathbb{Q}$.

\begin{answer}
Yes. For example
\[
S=\{q\in\mathbb{Q}:q<1\}
\]
has supremum $1$ in $\mathbb{Q}$ but no maximum.
\end{answer}

\item Can a set have a maximum without having a supremum?

\begin{answer}
No. If $m=\max S$, then $m$ is an upper bound of $S$, and since $m\in S$, every upper bound $u$ of $S$ satisfies $m\le u$. Hence $m$ is the least upper bound, so
\[
\sup S=m.
\]
This holds in any ordered set, with no completeness assumption needed.
\end{answer}

\item If $A$ is the lower half of a Dedekind cut, does $A$ have an upper bound in $\mathbb{Q}$?
Prove your answer directly from the cut axioms.

\begin{answer}
Yes. Choose any $b\in B$. Since every $a\in A$ satisfies $a<b$, the number $b$ is a rational upper bound of $A$.
\end{answer}

\item If $A\subseteq\mathbb{Q}$ is bounded above, must $\sup_{\mathbb{Q}}A$ exist?

\begin{answer}
No. For example,
\[
\{q\in\mathbb{Q}:q^2<2,\ q>0\}
\]
is bounded above in $\mathbb{Q}$ but has no rational supremum.
\end{answer}

\item Why is the notation
\[
\sup A
\]
potentially ambiguous in a foundational discussion?

\begin{answer}
The ambient ordered set is omitted. A set may have a supremum in $\mathbb{R}$ but not in $\mathbb{Q}$.
\end{answer}

\item Explain the difference between
\[
\sup_{\mathbb{Q}}A
\qquad\text{and}\qquad
\sup_{\mathbb{R}}A.
\]

\begin{answer}
$\sup_{\mathbb{Q}}A$ must itself be rational, while $\sup_{\mathbb{R}}A$ may be irrational.
\end{answer}

\item Consider
\[
A=\{q\in\mathbb{Q}:q<0\text{ or }q^2<2\}.
\]
Does $A$ have a supremum in $\mathbb{Q}$?

\begin{answer}
No. Its least upper bound would have to be $\sqrt2$, which is not rational.
\end{answer}

\item If $\mathbb{R}$ has already been constructed, what is $\sup_{\mathbb{R}}A$ for the set above?

\begin{answer}
\[
\sup_{\mathbb{R}}A=\sqrt2.
\]
\end{answer}

\item Why would using $\sup_{\mathbb{R}}A$ during the construction of $\mathbb{R}$ risk circular reasoning?

\begin{answer}
If $\mathbb{R}$ is being constructed from cuts, using completeness of $\mathbb{R}$ to define or justify the cuts assumes the object being constructed.
\end{answer}

\item Does ``$A$ has no greatest element'' imply ``$A$ has no supremum''?

\begin{answer}
No. For example,
\[
A=\{q\in\mathbb{Q}:q<1\}
\]
has no greatest element but has supremum $1$ in $\mathbb{Q}$.
\end{answer}

\item Does ``$A$ has no greatest element'' imply that $A$ is unbounded above?

\begin{answer}
No. The same set is bounded above by $1$.
\end{answer}

\end{enumerate}

\section{What Can Be Proved Using Only \texorpdfstring{$\mathbb{Q}$}{Q}?}

\begin{enumerate}[label=\arabic*., itemsep=10pt]

\item Before $\mathbb{R}$ is constructed, which number system are the elements of a Dedekind cut taken from?

\begin{answer}
Every element of a cut is rational. The construction begins entirely inside $\mathbb{Q}$.
\end{answer}

\item Which of the following can be proved using only the cut axioms and the order properties of $\mathbb{Q}$?
\begin{enumerate}[label=(\alph*)]
\item $A$ is bounded above.
\item $A$ has a rational supremum.
\item $A$ has no greatest element.
\item $B$ is upward closed.
\item Every element of $B$ is an upper bound of $A$.
\end{enumerate}

\begin{answer}
(a), (c), (d), and (e) can be proved. Statement (b) is false in general because $\mathbb{Q}$ is not complete.
\end{answer}

\item Why is it legitimate to use density of $\mathbb{Q}$ when constructing $\mathbb{R}$ from cuts?

\begin{answer}
Density of $\mathbb{Q}$ is already a theorem about rational numbers, so it does not depend on the existence of $\mathbb{R}$.
\end{answer}

\item Why is it not legitimate to use completeness of $\mathbb{R}$ at that stage?

\begin{answer}
Completeness of $\mathbb{R}$ is one of the properties the construction is intended to establish, so using it beforehand would be circular.
\end{answer}

\item Does the Archimedean property belong only to $\mathbb{R}$?

\begin{answer}
No. An Archimedean property can already be proved for $\mathbb{Q}$.
\end{answer}

\item State an Archimedean property that can be proved already for $\mathbb{Q}$.

\begin{answer}
One version is:
for every $x\in\mathbb{Q}$ there exists $n\in\mathbb{N}$ such that
\[
n>x.
\]
\end{answer}

\item If a proof about cuts invokes the Archimedean property, what must you check to make sure the proof is not circular?

\begin{answer}
You must check that the version used was proved from properties of $\mathbb{Q}$ or $\mathbb{N}$, rather than imported from completeness of $\mathbb{R}$.
\end{answer}

\item Suppose someone says:
``Since $A$ is bounded above, it has a least upper bound.''
What hidden assumption has been made?

\begin{answer}
The hidden assumption is completeness of the ambient ordered set.
\end{answer}

\item Suppose someone says:
``The cut has a boundary point $c$, and therefore $A=\{q\in\mathbb{Q}:q<c\}$.''
What has potentially been assumed without proof?

\begin{answer}
The argument assumes that such a boundary object $c$ already exists in some larger ordered set.
\end{answer}

\item Give an example of a statement about cuts that is true after $\mathbb{R}$ is constructed but cannot be used as a premise in constructing $\mathbb{R}$.

\begin{answer}
For example: ``Every nonempty bounded subset of $\mathbb{R}$ has a supremum.'' This is valid only after $\mathbb{R}$ and its completeness have been established.
\end{answer}

\end{enumerate}

\section{Rational Cuts and Irrational Cuts}

\begin{enumerate}[label=\arabic*., itemsep=10pt]

\item For $r\in\mathbb{Q}$, define the lower set
\[
A_r=\{q\in\mathbb{Q}:q<r\}.
\]

\begin{answer}
\[
A_r=\{q\in\mathbb{Q}:q<r\}.
\]
\end{answer}

\item What is the corresponding upper set $B_r$?

\begin{answer}
\[
B_r=\{q\in\mathbb{Q}:q\ge r\}.
\]
\end{answer}

\item Does $A_r$ have a greatest element?

\begin{answer}
No. Between any $q<r$ and $r$ there is another rational.
\end{answer}

\item Does $B_r$ have a least element?

\begin{answer}
Yes. Its least element is $r$.
\end{answer}

\item Prove that if a cut is produced by a rational number $r$ under the usual convention, then $B$ has a least element.

\begin{answer}
By definition,
\[
B_r=\{q\in\mathbb{Q}:q\ge r\},
\]
so $r\in B_r$ and $r\le b$ for every $b\in B_r$.
\end{answer}

\item Suppose $B$ has a least element $b$. Show that the cut is exactly the rational cut determined by $b$.

\begin{answer}
If $b=\min B$, then every rational $q<b$ cannot lie in $B$, so $q\in A$. Also no rational $q\ge b$ lies in $A$. Hence
\[
A=\{q:q<b\}.
\]
\end{answer}

\item Deduce the intrinsic characterization
\[
\text{the cut comes from a rational number}
\iff
B\text{ has a least element}.
\]

\begin{answer}
The two implications above prove
\[
B\text{ has a least element}
\iff
\text{the cut is induced by some }r\in\mathbb{Q}.
\]
\end{answer}

\item How would you characterize a ``new'' cut, one not corresponding to any rational number, using only the internal structure of $B$?

\begin{answer}
It is a cut whose upper set $B$ has no least element.
\end{answer}

\item Why is this characterization preferable, during the construction, to saying
``the boundary is irrational''?

\begin{answer}
Because ``irrational boundary'' already presupposes the existence of irrational real numbers. The property of $B$ is internal to $\mathbb{Q}$.
\end{answer}

\item For the cut associated with $\sqrt{2}$, why is there no least rational element of $B$?

\begin{answer}
If $b\in B$, then $b>\sqrt2$. By density, there exists rational $b'$ with
\[
\sqrt2<b'<b,
\]
so $b$ cannot be least.
\end{answer}

\end{enumerate}

\section{The Dangerous Word ``Boundary''}

\begin{enumerate}[label=\arabic*., itemsep=10pt]

\item What does the expression
\[
A=\{q\in\mathbb{Q}:q<c\}
\]
presuppose about $c$?

\begin{answer}
It presupposes that $c$ is already an element of some ordered set in which comparison with rationals is meaningful.
\end{answer}

\item If $c\in\mathbb{Q}$, is there any foundational problem with this definition?

\begin{answer}
No. If $c\in\mathbb{Q}$, the definition uses only rational order.
\end{answer}

\item If $c\notin\mathbb{Q}$ and $\mathbb{R}$ has not yet been constructed, what exactly is $c$ supposed to be?

\begin{answer}
Without an already constructed extension of $\mathbb{Q}$, such a $c$ has not yet been defined.
\end{answer}

\item Why can the geometric picture of a cut as ``all rationals to the left of a boundary'' be conceptually useful but logically dangerous?

\begin{answer}
The picture suggests a pre-existing endpoint, whereas in the construction the cut itself is the new number.
\end{answer}

\item In the Dedekind construction, is the real number first given and then encoded by a cut, or is the cut itself taken to be the real number?

\begin{answer}
The cut itself is taken to be the real number.
\end{answer}

\item Rewrite the statement
``this cut has boundary $\sqrt{2}$''
in a way that does not presuppose that $\sqrt{2}$ already exists as a real number.

\begin{answer}
One can say: define the object called $\sqrt2$ to be the cut
\[
\{q\in\mathbb{Q}:q<0\text{ or }q^2<2\}.
\]
\end{answer}

\item Suppose a cut has no rational boundary. Does that mean it has no meaning or no location in the order?

\begin{answer}
No. The cut itself specifies a position in the ordering even when no rational lies exactly at that position.
\end{answer}

\item What does it mean to say that a cut identifies a ``gap'' in $\mathbb{Q}$?

\begin{answer}
It means the partition behaves like a missing order-location in $\mathbb{Q}$: all rationals on one side lie below all rationals on the other, but no rational realizes the boundary.
\end{answer}

\end{enumerate}

\section{How an Irrational Number Is Defined Without Naming It First}

\begin{enumerate}[label=\arabic*., itemsep=10pt]

\item Consider
\[
A=\{q\in\mathbb{Q}:q<0\text{ or }q^2<2\}.
\]
Why can this set be defined without assuming the existence of the real number $\sqrt{2}$?

\begin{answer}
The defining condition uses only rational inequalities and rational squaring. No prior real number is required.
\end{answer}

\item Verify that this set is nonempty.

\begin{answer}
For example, $0\in A$.
\end{answer}

\item Verify that this set is a proper subset of $\mathbb{Q}$.

\begin{answer}
For example, $2\notin A$, so $A\neq\mathbb{Q}$.
\end{answer}

\item Prove that it is downward closed.

\begin{answer}
If $x\in A$ and $y<x$, then either $x<0$, which gives $y<0$, or $x\ge0$ and $x^2<2$. In the latter case, negative $y$ is automatically in $A$, while $0\le y<x$ gives $y^2<x^2<2$.
\end{answer}

\item Prove that it has no greatest element.

\begin{answer}
For negative $q$, choose a slightly larger negative rational. For nonnegative $q$ with $q^2<2$, one can explicitly choose a sufficiently small positive rational increment that keeps the square below $2$.
\end{answer}

\item Which part of that proof requires the most care?

\begin{answer}
The last step: constructing a larger rational while preserving $q^2<2$ without using $\sqrt2$.
\end{answer}

\item Once this set is proved to be a Dedekind cut, what object is taken to be $\sqrt{2}$ in the Dedekind construction?

\begin{answer}
The cut itself is defined to be the real number $\sqrt2$.
\end{answer}

\item Does every irrational real number need to have a simple algebraic description like $q^2<2$?

\begin{answer}
No. Most irrational numbers do not have a simple algebraic formula.
\end{answer}

\item If not, how can the construction still contain all real numbers?

\begin{answer}
The collection of all Dedekind cuts is taken as the set of real numbers. A cut need not have a simple finite description to exist as a set.
\end{answer}

\item What is conceptually wrong with asking:
``How do we define the cut belonging to an arbitrary irrational number before irrational numbers have been defined?''

\begin{answer}
It presupposes that the irrational number exists independently before the construction. The direction is reversed: the cut defines the number.
\end{answer}

\item After $\mathbb{R}$ has been constructed, how can one show that every real number corresponds to a Dedekind cut of $\mathbb{Q}$?

\begin{answer}
Given $\alpha\in\mathbb{R}$, define
\[
A_\alpha=\{q\in\mathbb{Q}:q<\alpha\}.
\]
One then verifies that this is a Dedekind cut.
\end{answer}

\end{enumerate}

\section{Density Versus Completeness}

\begin{enumerate}[label=\arabic*., itemsep=10pt]

\item State the density property of $\mathbb{Q}$.

\begin{answer}
For any rationals $a<b$, there exists rational $q$ with
\[
a<q<b.
\]
\end{answer}

\item Does density mean that every bounded subset of $\mathbb{Q}$ has a supremum in $\mathbb{Q}$?

\begin{answer}
No. Density concerns points between two existing points; completeness concerns least upper bounds.
\end{answer}

\item Give a bounded subset of $\mathbb{Q}$ with no supremum in $\mathbb{Q}$.

\begin{answer}
For example,
\[
\{q\in\mathbb{Q}:q^2<2,\ q>0\}.
\]
\end{answer}

\item Why does the existence of rationals arbitrarily close to a gap not fill the gap?

\begin{answer}
Density produces better approximations but does not produce the missing limiting point itself.
\end{answer}

\item Explain the distinction:
\[
\text{there are points arbitrarily close}
\qquad\text{versus}\qquad
\text{the limiting boundary belongs to the space}.
\]

\begin{answer}
``Arbitrarily close'' is an approximation statement. ``Belongs to the space'' is an existence statement.
\end{answer}

\item Can an ordered set be dense but incomplete?

\begin{answer}
Yes. $\mathbb{Q}$ is the standard example.
\end{answer}

\item Is $\mathbb{Q}$ dense in itself?

\begin{answer}
Yes.
\end{answer}

\item Is $\mathbb{Q}$ complete?

\begin{answer}
No.
\end{answer}

\item Why is there no contradiction between the previous two answers?

\begin{answer}
Density and completeness are different properties. One does not imply the other.
\end{answer}

\item Does the statement
``between any two rationals there is another rational''
say anything by itself about the existence of suprema?

\begin{answer}
No. It says nothing about whether bounded subsets have least upper bounds.
\end{answer}

\item In what sense is completeness a global property while density is a local order property?

\begin{answer}
Density is local because it concerns pairs of points. Completeness is global because it concerns arbitrary bounded subsets.
\end{answer}

\item Why are Dedekind cuts specifically designed to repair incompleteness rather than lack of density?

\begin{answer}
$\mathbb{Q}$ already has density. What it lacks are certain least upper bounds, and cuts are designed to add those missing order-locations.
\end{answer}

\end{enumerate}

\section{Consequences of ``No Greatest Element''}

\begin{enumerate}[label=\arabic*., itemsep=10pt]

\item Write again the quantified form of ``$A$ has no greatest element.''

\begin{answer}
\[
\forall a\in A\;\exists a'\in A\text{ such that }a<a'.
\]
\end{answer}

\item Prove that a nonempty set with no greatest element must be infinite.

\begin{answer}
Every finite nonempty linearly ordered set has a greatest element. Therefore a nonempty set with no greatest element must be infinite.
\end{answer}

\item Does no greatest element imply unboundedness above? Give a counterexample if not.

\begin{answer}
No. For example,
\[
\{q\in\mathbb{Q}:q<1\}
\]
has no greatest element but is bounded above.
\end{answer}

\item Starting from any $a_1\in A$, explain how to construct
\[
a_1<a_2<a_3<\cdots
\]
with every $a_n\in A$.

\begin{answer}
Choose $a_2>a_1$ in $A$, then $a_3>a_2$, and continue recursively.
\end{answer}

\item Does the existence of such a strictly increasing sequence imply that it converges?

\begin{answer}
No. A strictly increasing sequence may diverge or fail to converge in the ambient space.
\end{answer}

\item Does it imply that it is bounded?

\begin{answer}
No in general. For example, $a_n=n$ is strictly increasing and unbounded.
\end{answer}

\item For a Dedekind lower set, why is every such sequence bounded above?

\begin{answer}
Any $b\in B$ is an upper bound for every element of $A$.
\end{answer}

\item Let $a\in A$ and let $\varepsilon>0$ be rational.
Using downward closure, density of $\mathbb{Q}$, and the no-greatest-element property, prove that there exists $a'\in A$ such that
\[
a<a'<a+\varepsilon.
\]

\begin{answer}
Choose $z\in A$ with $a<z$. By density choose rational
\[
a<a'<\min(z,a+\varepsilon).
\]
Since $a'<z$ and $A$ is downward closed, $a'\in A$.
\end{answer}

\item Which hypotheses were actually needed in the previous proof?

\begin{answer}
No greatest element, downward closure, and density of $\mathbb{Q}$ were used.
\end{answer}

\item Does ``no greatest element'' alone imply that points of $A$ occur arbitrarily close above every $a\in A$ in an arbitrary ordered set?

\begin{answer}
No. In an ordered set with jumps, there may be a next element above $a$ and no points arbitrarily close to it.
\end{answer}

\item Why does the structure of $\mathbb{Q}$ matter?

\begin{answer}
Density of $\mathbb{Q}$ allows one to refine any strict inequality by inserting another rational.
\end{answer}

\item After $\mathbb{R}$ has been constructed, suppose $s=\sup_{\mathbb{R}}A$.
Why can $s$ not belong to $A$?

\begin{answer}
If $s\in A$, then no-greatest-element gives some $a\in A$ with $s<a$, contradicting that $s$ is an upper bound.
\end{answer}

\item Is the previous argument available before $\mathbb{R}$ is constructed?

\begin{answer}
Not if $s$ is being taken as a real supremum before $\mathbb{R}$ has been constructed.
\end{answer}

\end{enumerate}

\section{Properties of the Upper Half \texorpdfstring{$B$}{B}}

\begin{enumerate}[label=\arabic*., itemsep=10pt]

\item Given a lower-set formulation of a cut, define
\[
B=\mathbb{Q}\setminus A.
\]

\begin{answer}
\[
B=\mathbb{Q}\setminus A.
\]
\end{answer}

\item Prove that $B$ is nonempty.

\begin{answer}
Since $A$ is proper, there exists a rational not in $A$, hence in $B$.
\end{answer}

\item Prove that $B$ is upward closed.

\begin{answer}
If $b\in B$ and $c>b$, then $c\notin A$; otherwise downward closure would force $b\in A$. Hence $c\in B$.
\end{answer}

\item Prove that every element of $B$ is greater than every element of $A$.

\begin{answer}
If some $b\in B$ satisfied $b\le a$ for $a\in A$, then downward closure of $A$ would force $b\in A$, contradiction.
\end{answer}

\item Prove that every element of $B$ is an upper bound of $A$.

\begin{answer}
Since every $a\in A$ satisfies $a<b$, each $b\in B$ is an upper bound of $A$.
\end{answer}

\item Is every upper bound of $A$ necessarily an element of $B$?

\begin{answer}
Yes, for rational upper bounds. If $u\in\mathbb{Q}$ is an upper bound of $A$, then $u\notin A$, because $A$ has no greatest element. Hence $u\in B$.
\end{answer}

\item Must $B$ have a least element?

\begin{answer}
No.
\end{answer}

\item When does $B$ have a least element?

\begin{answer}
Exactly when the cut is induced by a rational number.
\end{answer}

\item Can $B$ have a greatest element?

\begin{answer}
No. Since $B$ is upward closed and nonempty, it contains arbitrarily large rationals.
\end{answer}

\item Can $B$ be finite?

\begin{answer}
No. The same reasoning shows it is infinite.
\end{answer}

\item Suppose $b\in B$ and $b'<b$.
Must $b'$ lie in $B$?

\begin{answer}
Not necessarily. It may lie in $A$ or $B$.
\end{answer}

\item Suppose $b\in B$ and $b'>b$.
Must $b'$ lie in $B$?

\begin{answer}
Yes. Upward closure gives $b'\in B$.
\end{answer}

\item Compare the asymmetry:
\[
A\text{ has no greatest element}
\]
while
\[
B\text{ may or may not have a least element}.
\]
Why is this asymmetry useful?

\begin{answer}
It makes rational cuts identifiable by the existence of $\min B$, while keeping the lower-set convention uniform for all cuts.
\end{answer}

\end{enumerate}

\section{Approximation Across the Cut}

\begin{enumerate}[label=\arabic*., itemsep=10pt]

\item Let $(A,B)$ be a Dedekind cut.
What would it mean for $A$ and $B$ to be separated by a positive rational gap?

\begin{answer}
It would mean that there exists $\delta>0$ such that
\[
b-a\ge\delta
\]
for every $a\in A$ and $b\in B$.
\end{answer}

\item Formulate precisely the statement
\[
\forall\varepsilon\in\mathbb{Q},\ \varepsilon>0,
\quad
\exists a\in A,\ \exists b\in B
\quad\text{such that}\quad
b-a<\varepsilon.
\]

\begin{answer}
\[
\forall\varepsilon\in\mathbb{Q}_{>0}\;
\exists a\in A\;\exists b\in B
\text{ such that }b-a<\varepsilon.
\]
\end{answer}

\item Why is this statement stronger than merely saying that every $a\in A$ is less than every $b\in B$?

\begin{answer}
Order separation only says $a<b$. The approximation statement says the distance can be made arbitrarily small.
\end{answer}

\item Try to prove the approximation statement using only properties of $\mathbb{Q}$ and the cut axioms.

\begin{answer}
One standard proof subdivides a rational interval from some $a_0\in A$ to some $b_0\in B$ into sufficiently many equal rational pieces and finds adjacent subdivision points lying on opposite sides of the cut.
\end{answer}

\item Where does the Archimedean property enter the proof?

\begin{answer}
It is used to choose enough subdivisions so that the mesh size is smaller than $\varepsilon$.
\end{answer}

\item Which version of the Archimedean property is sufficient?

\begin{answer}
It suffices that for every positive rational $\varepsilon$ and positive rational $M$, there exists $n\in\mathbb{N}$ with
\[
M/n<\varepsilon.
\]
\end{answer}

\item Why does invoking that version not require completeness of $\mathbb{R}$?

\begin{answer}
This is a property of $\mathbb{Q}$ and $\mathbb{N}$ and can be proved without real completeness.
\end{answer}

\item Could the approximation result be proved directly from density of $\mathbb{Q}$ alone?

\begin{answer}
Not by density alone in the most direct argument. Density inserts points between known rationals but does not by itself give a uniform finite subdivision argument.
\end{answer}

\item If there were a fixed rational $\delta>0$ such that
\[
b-a\ge\delta
\]
for every $a\in A$ and $b\in B$, what would go wrong?

\begin{answer}
Such a gap would allow rational points between the two halves that belong to neither side, contradicting that $A\cup B=\mathbb{Q}$, once sufficiently fine rational subdivisions are considered.
\end{answer}

\item How does the no-greatest-element property help rule out such a gap?

\begin{answer}
If $a\in A$, there are larger elements of $A$. Thus $A$ cannot terminate a positive distance before all of $B$.
\end{answer}

\item Does the approximation property imply that there exists a rational number lying exactly ``between'' $A$ and $B$?

\begin{answer}
No. Arbitrarily small gaps do not imply that a rational realizes the boundary.
\end{answer}

\item Why is the difference between ``arbitrarily close'' and ``attained boundary'' important here?

\begin{answer}
A sequence of rational approximations may approach a missing location without any rational being equal to that location.
\end{answer}

\end{enumerate}

\section{Arithmetic on Dedekind Cuts}

\begin{enumerate}[label=\arabic*., itemsep=10pt]

\item If real numbers are Dedekind cuts, what should it mean to add two real numbers?

\begin{answer}
One defines an operation on cuts that corresponds to rational addition and then proves that the result is again a cut.
\end{answer}

\item Given cuts $A$ and $C$, suggest a definition for their sum using rational sums $a+c$.

\begin{answer}
A natural definition is
\[
A+C=\{q\in\mathbb{Q}:\exists a\in A,\exists c\in C,\ q<a+c\},
\]
or an equivalent downward-closed formulation.
\end{answer}

\item Why is it not enough merely to define a set? What must be proved about the resulting set?

\begin{answer}
Because a real number is a cut, the output must satisfy all Dedekind-cut axioms.
\end{answer}

\item Which cut axioms must be checked for addition?

\begin{answer}
Nonemptiness, properness, downward closure, and absence of a greatest element.
\end{answer}

\item Why is multiplication more complicated than addition?

\begin{answer}
Products interact with order differently depending on signs.
\end{answer}

\item Why do signs create complications when defining multiplication of cuts?

\begin{answer}
Multiplication by a negative number reverses inequalities, so a single naive formula does not work uniformly.
\end{answer}

\item What should the additive identity look like as a cut?

\begin{answer}
The additive identity is the cut
\[
\{q\in\mathbb{Q}:q<0\}.
\]
\end{answer}

\item What should the cut representing a rational number $r$ look like?

\begin{answer}
\[
A_r=\{q\in\mathbb{Q}:q<r\}.
\]
\end{answer}

\item What should negation of a cut accomplish order-theoretically?

\begin{answer}
It should produce a cut whose represented value is the order-theoretic reflection of the original one about $0$.
\end{answer}

\item Why is defining the reciprocal of a positive cut more delicate than defining its sum with another cut?

\begin{answer}
Reciprocals reverse inequalities on positive numbers and require care near $0$.
\end{answer}

\item After defining arithmetic, what must be proved before one may call the resulting structure an ordered field?

\begin{answer}
One must verify the field axioms, compatibility with the order, and that the rational embedding preserves the arithmetic.
\end{answer}

\end{enumerate}

\section{Completeness After the Construction}

\begin{enumerate}[label=\arabic*., itemsep=10pt]

\item What is the logical order of the Dedekind construction?
Arrange the following:
\[
\mathbb{Q},\qquad
\text{Dedekind cuts},\qquad
\mathbb{R},\qquad
\text{least-upper-bound property}.
\]

\begin{answer}
\[
\mathbb{Q}
\longrightarrow
\text{Dedekind cuts}
\longrightarrow
\mathbb{R}
\longrightarrow
\text{proof of the least-upper-bound property}.
\]
\end{answer}

\item Why must completeness be proved after the real numbers have been defined?

\begin{answer}
Completeness is a theorem about the newly defined structure, not a premise that may be assumed beforehand.
\end{answer}

\item Suppose $\mathcal{S}$ is a nonempty set of Dedekind cuts bounded above.
What kind of object should $\sup \mathcal{S}$ itself be?

\begin{answer}
It should itself be a Dedekind cut, hence an element of the constructed $\mathbb{R}$.
\end{answer}

\item A common construction defines the supremum of a family of cuts by a union.
Why is
\[
\bigcup_{A\in\mathcal{S}} A
\]
a natural candidate?

\begin{answer}
The union collects every rational lying below at least one cut in the family, which is exactly what a least upper bound should do.
\end{answer}

\item Which Dedekind-cut axioms must be verified for that union?

\begin{answer}
One must verify nonemptiness, properness, downward closure, and no greatest element.
\end{answer}

\item Why does this proof of completeness avoid circular reasoning?

\begin{answer}
The candidate supremum is constructed explicitly from sets of rationals, without invoking a previously existing supremum in $\mathbb{R}$.
\end{answer}

\item What is the conceptual difference between:
\[
\text{assuming every bounded set has a supremum}
\]
and
\[
\text{constructing the supremum explicitly as another cut}?
\]

\begin{answer}
The first is an existence axiom or theorem. The second provides a concrete object and proves that it has the required property.
\end{answer}

\item Once completeness of the cut construction has been proved, why is it thereafter legitimate to write
\[
\sup_{\mathbb{R}} S
\]
for suitable bounded sets $S\subseteq\mathbb{R}$?

\begin{answer}
Because completeness has now been established as a theorem of the constructed real number system.
\end{answer}

\item At that later stage, how can one reinterpret a Dedekind cut geometrically as
\[
A=\{q\in\mathbb{Q}:q<\alpha\}
\]
for some $\alpha\in\mathbb{R}$?

\begin{answer}
For $\alpha\in\mathbb{R}$, its corresponding lower rational set is
\[
\{q\in\mathbb{Q}:q<\alpha\}.
\]
\end{answer}

\item Why would using that geometric representation at the beginning of the construction reverse the logical dependency?

\begin{answer}
It would assume the real number $\alpha$ exists before the construction that is supposed to define it.
\end{answer}

\end{enumerate}

\section{Mixed Assumption Checks}

For each statement, decide whether it is true, false, or meaningless without further specification.
If it is ambiguous, identify the hidden assumption.

\begin{enumerate}[label=\arabic*., itemsep=10pt]

\item Every Dedekind cut has a supremum.

\begin{answer}
Ambiguous. It depends on the ambient ordered set. In $\mathbb{R}$, yes; in $\mathbb{Q}$, not always.
\end{answer}

\item Every Dedekind lower set is bounded above.

\begin{answer}
True. Any $b\in B$ is a rational upper bound of $A$.
\end{answer}

\item Every Dedekind lower set has no maximum.

\begin{answer}
True. This is part of the standard definition.
\end{answer}

\item Every Dedekind upper set has a minimum.

\begin{answer}
False. It has a minimum exactly for rational cuts under the usual convention.
\end{answer}

\item If $A$ has no greatest element, then $A$ is unbounded above.

\begin{answer}
False.
\end{answer}

\item If $A$ is bounded above, then $\sup A$ exists.

\begin{answer}
False in general. It requires completeness of the ambient ordered set.
\end{answer}

\item If $A$ is a Dedekind lower set, then $\sup_{\mathbb{Q}}A$ exists.

\begin{answer}
False.
\end{answer}

\item If $A$ is a Dedekind lower set, then $\sup_{\mathbb{R}}A$ exists.

\begin{answer}
True after $\mathbb{R}$ has been constructed and proved complete.
\end{answer}

\item Every cut has a boundary point.

\begin{answer}
Ambiguous. Geometrically yes after embedding into $\mathbb{R}$; not necessarily as a rational point.
\end{answer}

\item Every cut has a rational boundary point.

\begin{answer}
False.
\end{answer}

\item Every cut corresponds to exactly one real number.

\begin{answer}
True after the Dedekind construction, because cuts are the real numbers in that model.
\end{answer}

\item Every real number is first defined independently and then represented by a cut.

\begin{answer}
False. In the construction, the cut itself defines the real number.
\end{answer}

\item The density of $\mathbb{Q}$ implies that $\mathbb{Q}$ has no gaps.

\begin{answer}
False. $\mathbb{Q}$ is dense but incomplete.
\end{answer}

\item A set can be bounded above and still fail to have a supremum in its ambient ordered set.

\begin{answer}
True.
\end{answer}

\item A set can have no greatest element and still have a supremum.

\begin{answer}
True.
\end{answer}

\item If $B$ has a least element, the cut corresponds to a rational number.

\begin{answer}
True.
\end{answer}

\item If $B$ has no least element, the cut corresponds to a new real number not represented by a rational.

\begin{answer}
True, meaning a cut not induced by any rational.
\end{answer}

\item The proof that cuts approximate each other arbitrarily closely must use completeness of $\mathbb{R}$.

\begin{answer}
False.
\end{answer}

\item The Archimedean property cannot be used before $\mathbb{R}$ is constructed.

\begin{answer}
False. Appropriate Archimedean properties can be proved already in $\mathbb{Q}$.
\end{answer}

\item Once $\mathbb{R}$ is constructed and proved complete, one may freely reinterpret cuts using real suprema.

\begin{answer}
True.
\end{answer}

\end{enumerate}

\section{Dependency Questions}

For each proposed argument, identify whether it is logically valid at the stage where real numbers are being constructed from rational Dedekind cuts.

\begin{enumerate}[label=\arabic*., itemsep=10pt]

\item ``$A$ is bounded above, therefore let $s=\sup A$.''

\begin{answer}
Invalid without specifying a complete ambient ordered set. Boundedness alone does not imply existence of a supremum in $\mathbb{Q}$.
\end{answer}

\item ``Choose $b\in B$. Then $b$ is an upper bound of $A$.''

\begin{answer}
Valid. The order-separation property of a cut makes every $b\in B$ an upper bound of $A$.
\end{answer}

\item ``Since $\mathbb{Q}$ is dense, choose a rational number strictly between two distinct rationals.''

\begin{answer}
Valid. Density of $\mathbb{Q}$ is available before constructing $\mathbb{R}$.
\end{answer}

\item ``The cut representing $\sqrt{2}$ consists of all rationals below the already existing real number $\sqrt{2}$.''

\begin{answer}
Invalid as a foundational definition. It presupposes the real number $\sqrt2$.
\end{answer}

\item ``Define the object called $\sqrt{2}$ to be the cut
\[
\{q\in\mathbb{Q}:q<0\text{ or }q^2<2\}.
\]
''

\begin{answer}
Valid, provided the displayed set is proved to satisfy the cut axioms.
\end{answer}

\item ``Because $\mathbb{R}$ is complete, the cut has a least upper bound.''

\begin{answer}
Invalid during the construction. It assumes the completeness of the structure being constructed.
\end{answer}

\item ``Because the Archimedean property holds in $\mathbb{Q}$, choose $n\in\mathbb{N}$ such that $n\varepsilon$ exceeds a prescribed rational number.''

\begin{answer}
Valid if that Archimedean statement has already been proved for $\mathbb{Q}$.
\end{answer}

\item ``Because $A$ has no greatest element, choose $a'\in A$ with $a<a'$.''

\begin{answer}
Valid. This is exactly the no-greatest-element axiom.
\end{answer}

\item ``Because $B$ is nonempty, choose an upper bound of $A$.''

\begin{answer}
Valid. Any $b\in B$ is an upper bound of $A$.
\end{answer}

\item ``Because every cut is itself defined to be a real number, no prior real boundary is needed.''

\begin{answer}
Valid. This is the conceptual basis of the construction.
\end{answer}

\end{enumerate}

\end{document}
